% Generated by tools/edn_latex.py from reports/edn.md.
% Do not edit: change reports/edn.md and run `make edn`.
\ednentry{
  id = {EDN-11},
  title = {Keep \texorpdfstring{\texttt{shap}}{shap} out of the API image; serving uses the booster's native SHAP export},
  date = {2026-09-22},
  milestone = {M4: Model Deployment},
  topic = {ML System Design, Explainability},
  participants = {@lukas2510},
  decision = {Amends EDN-09's serving-path clause. FR-08's per-request explanation is computed from the trained booster's own SHAP export (LightGBM \texttt{predict(pred\_\allowbreak{}contrib=\allowbreak{}True)}, or CatBoost \texttt{get\_\allowbreak{}feature\_\allowbreak{}importance(type=\allowbreak{}\textquotedbl{}ShapValues\textquotedbl{})} if CatBoost is ever the deployed model), not the \texttt{shap} package. \texttt{shap} stays a training/notebook-only dependency, used for offline global analysis (summary plots, dataset-level importance), and is never installed in the API image. The Python 3.12 bump from EDN-09 is unaffected and still stands, since it is what lets \texttt{shap} run at all for that offline use.},
  rationale = {\emph{Alternatives considered:}
\begin{itemize}[leftmargin=*, topsep=2pt]
\item \textbf{Option A: keep \texttt{shap.\allowbreak{}TreeExplainer} in the serving path, as EDN-09 originally set FR-08 to.}
\newline Pros: one implementation for both the offline/notebook analysis and the serving path, nothing to keep in sync.
\newline Cons: \texttt{shap} unconditionally imports \texttt{numba} and \texttt{llvmlite} at module load, even to use only \texttt{TreeExplainer}, confirmed by deleting them and watching a bare \texttt{import shap} fail; a real Docker build of the serving stack (\texttt{shap}, \texttt{lightgbm}, \texttt{mapie}, \texttt{fastapi}, \texttt{uvicorn}, \texttt{pydantic}, \texttt{numpy}, \texttt{pandas}) measured 702 MB, leaving only about 30\,\% of NFR-04's 1 GB budget for the baked-in model artefact, app code and everything else still to be added.
\item \textbf{Option B (chosen): compute serving-time explanations from the booster's native SHAP export; keep \texttt{shap} for offline analysis only.}
\newline Pros: the same Docker build without \texttt{shap} measured 486 MB, about 51\,\% headroom instead of 30\,\%; the native export is verified bit-identical to \texttt{shap.\allowbreak{}TreeExplainer} (max absolute difference 0.0 across all features and the base value in a direct comparison); the pattern generalises to CatBoost, the brief's stated challenger model (EDN-02), via its own native SHAP support, so it does not need revisiting if the challenger wins; \texttt{shap} is still fully available where it is actually needed, offline global analysis, unconstrained by image size.
\newline Cons: two call sites compute the same values (the booster's native export in the API, \texttt{shap.\allowbreak{}TreeExplainer} in training/notebook code) instead of one; a training-time test is needed to keep them provably in sync (added to FR-08's ``Verified by'').
\end{itemize}
\par
\emph{Why:} The image-size argument is independent of the Python-version problem EDN-09 solved, and EDN-09 folded both into one decision without weighing the size evidence, which was not yet measured at the time. Measured directly (Docker builds, a deletion test proving \texttt{numba}/\texttt{llvmlite} are unconditional runtime imports, and a numeric comparison proving the native export and \texttt{shap.\allowbreak{}TreeExplainer} produce identical values), option B costs nothing in correctness and gives back real budget headroom under NFR-04, while also removing a dependency choice tied to a single model family.},
  ai = {Information seeking, Alternative assessment, Recommendation},
  response = {Accepted},
  assessment = {AI had proposed the native-export approach earlier in the PR \#19 review, before EDN-09 was decided. Asked whether it was still relevant after the Python 3.12 bump, AI separated the two motivations (version conflict, now resolved; image size, still open), then built real Docker images for both options and a deletion test to confirm \texttt{numba}/\texttt{llvmlite} are unconditional, rather than assuming it from documentation. Asked whether this was the best option or if there were better ones, AI additionally ruled out a leaner Docker build (does not work, proven by the deletion test) and dropping per-instance SHAP entirely (would weaken FR-08 and lose the course's stated M3 credit for using SHAP), and pointed out CatBoost's equivalent native support. Lukas reviewed the evidence and confirmed option B.},
  aievidence = {Claude Code session on 2026-09-22 while reviewing PR \#19: AI's original \texttt{pred\_\allowbreak{}contrib} suggestion predated the Python 3.12 bump; Lukas asked ``is this still relevant?'', AI re-derived it with concrete measurements; Lukas asked ``is this what we would want? or are there better options?'', AI built real Docker images (702 MB vs. 486 MB) and ruled out the alternatives; Lukas replied ``ok apply it.''},
  evidence = {\href{https://github.com/mlops-2627q1-mds-upc/MLOPS_Recommenditos/blob/main/docs/docs/specification.md}{Specification} FR-08; \href{https://github.com/mlops-2627q1-mds-upc/MLOPS_Recommenditos/blob/main/docs/docs/project-brief.md}{project brief} section 6 (tooling constraints); \hyperref[EDN-09]{EDN-09}; \href{https://github.com/mlops-2627q1-mds-upc/MLOPS_Recommenditos/pull/19}{PR \#19}.},
}
